\section{Repository-State and Agent-State Scalability}
\label{sec:scalability}

The paper retains a distinction between repository-state scalability and agent-state scalability, but the latter now has an explicit definition.

\emph{Repository-state scalability} concerns storing, serving, and reconstructing repository objects and history as repository count and size grow. Cursor Continuity is a contemporary example of durable object/log state with replaceable repository-serving compute \citep{marti2026continuity}.

\emph{Agent-state scalability} is an architectural objective: the ability to increase the number or duration of concurrent long-horizon agentic programmes without proportionally increasing persistent \emph{high-capability reasoning-session state} requirements.

Figure~\ref{fig:scalability-stack} is conceptual.

\begin{figure}[htbp]
\centering
\begin{tikzpicture}[
node distance=4mm,
repo/.style={draw,rounded corners,align=center,minimum width=8cm,minimum height=8mm},
agent/.style={draw,dashed,rounded corners,align=center,minimum width=8cm,minimum height=8mm},
arrow/.style={-{Latex[length=1.8mm]},thick}
]
\node[repo] (object) {Durable object / log state};
\node[repo,below=of object] (serve) {Replaceable repository-serving compute};
\node[repo,below=of serve] (eng) {Durable engineering repository state};
\node[agent,below=of eng] (reason) {Replaceable high-capability reasoning};
\node[agent,below=of reason] (exec) {Replaceable constrained execution workers};
\draw[arrow] (object)--(serve);
\draw[arrow] (serve)--(eng);
\draw[arrow] (eng)--(reason);
\draw[arrow] (reason)--(exec);
\node[right=5mm of serve,align=left,font=\small] {repository-state\\scalability};
\node[right=5mm of reason,align=left,font=\small] {agent-state\\scalability};
\end{tikzpicture}

\caption{CONCEPTUAL --- NOT EMPIRICAL DATA. Repository-serving durability and high-capability reasoning-session persistence are distinct layers. Cursor Continuity supports only the upper repository-infrastructure precedent.}
\label{fig:scalability-stack}
\end{figure}

The two problems are independent. A scalable Git service does not imply that agents can resume without conversational history. Conversely, reasoning-session disposal does not require novel Git infrastructure.

Long-context inference creates serving-state and processing requirements, but production serving is more complex than one conversation occupying one dedicated GPU. PagedAttention manages KV-cache allocation \citep{kwon2023vllm}; DistServe separates prefill and decode resources \citep{zhong2024distserve}; Mooncake uses disaggregated KV-cache infrastructure \citep{qin2024mooncake}. Consequently, terminating a reasoning session does not map cleanly to a provider-internal resource saving.

The only defensible infrastructure statement at v0.1 is conditional: if a workload can reduce persistent high-capability reasoning-state requirements while preserving useful outcomes and without paying the same cost in reconstruction, a serving system \emph{may} have more scheduling flexibility or useful concurrency. This is an \textbf{IMPLICATION}, not a measured result.

A direct agent-state scalability claim therefore requires future workload-level serving experiments, not token accounting alone.
